\documentclass[11pt,aspectratio=169]{beamer}

\usetheme{default}
\setbeamertemplate{navigation symbols}{}

\usepackage[T1]{fontenc}
\usepackage[utf8]{inputenc}
\usepackage{xcolor}
\usepackage{listings}

% Listings (Java)
\lstset{
	language=Java,
	basicstyle=\ttfamily\footnotesize,
	keywordstyle=\bfseries\color{blue!70!black},
	commentstyle=\itshape\color{teal!60!black},
	stringstyle=\color{red!60!black},
	showstringspaces=false,
	columns=flexible,
	tabsize=2,
	breaklines=true,
	frame=single,
	numbers=none,
	numbersep=6pt
}

\title{Inheritance vs Composition in \textbf{OOP}}
\author{Micro‑Course (10 minutes)}
\date{}

\begin{document}
	
	\begin{frame}
		\titlepage
	\end{frame}
	
	% Slide 1 — OOP in 60 seconds
	\begin{frame}[fragile]{\textbf{OOP} in 60 seconds}
		\begin{itemize}
			\item We model software as \textbf{objects}: \textbf{state} + \textbf{behavior}.
			\item A \textbf{class} is a blueprint; it defines attributes and methods.
			\item Why it matters: \textbf{encapsulation}, \textbf{reuse}, \textbf{polymorphism}.
		\end{itemize}
		\medskip
		\begin{lstlisting}
			// Object = state + behavior
			final class BankAccount {
				private final String owner;            // state
				private int balance;                   // state
				public BankAccount(String o, int b){ owner=o; balance=b; }
				public void deposit(int amount){ if(amount>0) balance += amount; } // behavior
			}
		\end{lstlisting}
	\end{frame}
	
	% Slide 2 — Inheritance (is-a)
	\begin{frame}[fragile]{\textbf{Inheritance} (\textbf{is-a})}
		\begin{itemize}
			\item A subclass \textbf{extends} a base class to reuse and specialize behavior.
			\item Use when subtype truly \textbf{is-a} base and respects \textbf{LSP} (Liskov Substitution Principle).
		\end{itemize}
		\medskip
		\begin{lstlisting}
			abstract class Shape { abstract int area(); }
			
			final class Rectangle extends Shape {
				private final int w, h;
				Rectangle(int w, int h){ this.w=w; this.h=h; }
				@Override int area(){ return w * h; }
			}
		\end{lstlisting}
		\textbf{Pros}: simple reuse, \textbf{polymorphism} via base type.
		\quad \textbf{Cons}: tight coupling, fragile base, LSP risks.
	\end{frame}
	
	% Slide 3 — Composition (has-a, delegation)
	\begin{frame}[fragile]{\textbf{Composition} (\textbf{has-a}, \textbf{delegation})}
		\begin{itemize}
			\item A class \textbf{has} components and \textbf{delegates} work to them.
			\item Use when behavior varies across instances or may change over time.
		\end{itemize}
		\medskip
		\begin{lstlisting}
			interface Engine { void start(); int hp(); }
			
			final class ElectricEngine implements Engine {
				public void start() { /* ... */ }
				public int hp() { return 250; }
			}
			
			final class Car {
				private final Engine engine;                  // has-a
				Car(Engine e){ this.engine = e; }
				void start(){ engine.start(); }               // delegation
				int power(){ return engine.hp(); }
			}
		\end{lstlisting}
		\textbf{Pros}: flexibility, testability, better \textbf{encapsulation}.
		\quad \textbf{Cons}: a bit more wiring.
	\end{frame}
	
	% Slide 4 — Side-by-side comparison
	\begin{frame}[fragile]{\textbf{Inheritance} vs \textbf{Composition}}
		\begin{columns}[T,onlytextwidth]
			\column{0.5\textwidth}
			\textbf{Inheritance (\textit{is-a})}
			\begin{itemize}
				\item \textbf{Pros}: shared code, natural \textbf{polymorphism}.
				\item \textbf{Cons}: tight coupling, \textbf{fragile base class}.
				\item \textbf{Use when}: stable base API; subtype truly \textbf{is-a} supertype and honors \textbf{LSP}.
			\end{itemize}
			
			\column{0.5\textwidth}
			\textbf{Composition (\textit{has-a})}
			\begin{itemize}
				\item \textbf{Pros}: replaceable parts, multiple change axes, strong \textbf{encapsulation}.
				\item \textbf{Cons}: more boilerplate/delegation.
				\item \textbf{Use when}: need runtime variability; avoid deep hierarchies.
			\end{itemize}
		\end{columns}
		\medskip
		\textbf{Rule of thumb}: Prefer \textbf{composition} for flexibility; use \textbf{inheritance} for true subtyping with stable behavior.
	\end{frame}
	
	% Slide 5 — Small refactor example (Decorator)
	\begin{frame}[fragile]{Adding behavior via \textbf{composition} (Decorator)}
		\begin{lstlisting}
			final class LoggingEngine implements Engine {
				private final Engine inner;                     // has-a
				LoggingEngine(Engine e){ this.inner = e; }
				public void start(){
					System.out.println("[LOG] start");            // add behavior
					inner.start();                                // delegate
				}
				public int hp(){ return inner.hp(); }
			}
		\end{lstlisting}
		\begin{itemize}
			\item No subclass of \textbf{Car}; no change to \textbf{Engine} types.
			\item We \textbf{compose} features instead of building fragile trees.
		\end{itemize}
	\end{frame}
	
	% Slide 6 — Takeaways + mini-exercise
	\begin{frame}[fragile]{\textbf{Takeaways} \& Mini‑exercise}
		\textbf{Takeaways}
		\begin{itemize}
			\item \textbf{Encapsulation}: hide internals; expose a small, safe API.
			\item \textbf{Inheritance} (\textbf{is-a}) vs \textbf{composition} (\textbf{has-a}): choose by \textbf{LSP}, variability, coupling.
			\item Favor \textbf{composition} when behavior may change or combine.
		\end{itemize}
		\medskip
		\textbf{30‑second exercise}
		\begin{itemize}
			\item Where would you add logging or caching using \textbf{composition} in the Car/Engine example?
			\item When would \textbf{inheritance} still be the cleaner choice?
		\end{itemize}
	\end{frame}
	
\end{document}